\section*{\heiti\zihao{-2}摘\qquad 要}

数学表达式识别与LaTeX生成是计算机视觉与自然语言处理交叉领域的重要研究课题，在文档数字化、教育辅助和学术研究等领域具有广泛应用前景。
传统方法多依赖卷积神经网络处理数学表达式图像，受限于局部感受野，难以建模数学符号间的长程结构依赖；而手写体表达式因笔迹风格差异大、结构层次复杂，识别难度进一步加剧。
为此，本文基于 PyTorch 实现了一种以 Vision Transformer（ViT）为编码器、标准 Transformer 为解码器的端到端模型框架，实现数学表达式图像的结构化识别及 LaTeX 序列生成，并在印刷体与手写体两类数据集上系统评估其性能。

模型方面，ViT 编码器将输入图像划分为固定尺寸的图像块，通过多层预归一化 Transformer 块的自注意力机制建模图像块间的全局依赖；解码器结合因果自注意力与交叉注意力机制，分别建模目标序列内部的上下文依赖以及与图像编码表示之间的对齐关系，输出层叠加基于训练语料词频的对数先验偏置。
训练方面，采用掩蔽交叉熵损失忽略填充位置、AdamW 优化器配合 Transformer 预热—衰减学习率调度与全局范数梯度裁剪、固定随机种子的验证集划分以及早停机制；推理阶段实现贪婪解码与带 GNMT 长度惩罚的束搜索解码。

实验在印刷体数据集 \texttt{im2latex-100k} 与手写体数据集 \texttt{MathWriting-human} 上以完全相同的架构与训练配方进行。
在印刷体测试集上，模型取得 $90.14\%$ 的词元级准确率，自回归生成的 BLEU-4 分数为 $0.6767$（贪婪解码）与 $0.7012$（束搜索），以 $7.73$M 的轻量参数规模与 $50\times200$ 的低分辨率输入验证了纯 Transformer 架构在该任务上的有效性。
训练稳定性分析提供了梯度裁剪必要性的直接实验证据：在大批量配合预热调度的设定下，无裁剪训练曾在学习率峰值附近发散并坍缩为退化解，引入裁剪后训练全程平稳收敛。

手写体对比实验则定量揭示了直接迁移印刷体配方的局限：测试集词元级准确率降至 $56.63\%$，BLEU-4 为 $0.0833$，且束搜索不再带来增益。
通过对验证—测试性能鸿沟、训练曲线分化与输入几何失真的分析，本文将该性能差距分解为书写者分布偏移、训练样本量不足与宽高比失真三个可分别施治的因素，并据此提出全量数据训练、数据增强与保宽高比预处理等改进方向。

研究结果表明，基于 Vision Transformer 的端到端模型在印刷体数学表达式识别任务上具有良好的性能表现；针对手写场景的对比实验在严格控制变量的前提下刻画了跨书写者泛化的核心挑战，为后续手写数学表达式识别研究提供了可复现的基线方法与明确的改进路径。

\vspace{0.5cm}
\noindent\textbf{关键词：}手写数学表达式识别；LaTeX 转换；Vision Transformer；注意力机制；深度学习；图像识别
